% ============================================================
%  Entorns didàctics locals del volum de categories.
%  Definits aquí (no a l'estil compartit) per no acoblar-hi
%  el projecte. Reutilitzen tcolorbox, etoolbox i la paleta
%  ja carregats per aprendes.sty en mode «sobri».
%
%    \begin{objectius} ... \end{objectius}   -> bloc d'obertura
%    \begin{resumcap}   ... \end{resumcap}   -> bloc de tancament
%
%  Dins d'aquests blocs s'hi poden usar els encapçalaments
%  \apartat{...} per separar prerequisits/objectius/etc.
% ============================================================

% Color propi per als blocs didàctics: verd pissarra derivat de la
% paleta existent (cexer = verd), enfosquit per contrastar amb el fons.
\colorlet{DidacticVerd}{cexer!55!black}

% Estils de bloc: mateixa gramàtica visual que blocenunciat/blocobservacio.
\tcbset{
  blocdidactic/.style={
    enhanced, breakable, sharp corners,
    frame hidden, boxrule=0pt,
    colback=DidacticVerd!6,
    borderline west={2.0pt}{0pt}{DidacticVerd},
    left=8pt, right=4pt, top=5pt, bottom=5pt,
    before skip=1.0em, after skip=1.0em,
  },
}

% Encapçalament intern d'un apartat del bloc (Prerequisits, Objectius...).
\newcommand{\apartat}[1]{\par\smallskip{\color{DidacticVerd}\ifpdftex\scshape\else\sffamily\bfseries\fi #1}\par\nobreak\smallskip}

% Bloc d'obertura del capítol.
\newenvironment{objectius}{%
  \begin{tcolorbox}[blocdidactic]%
  \ignorespaces
}{\end{tcolorbox}}

% Bloc de tancament del capítol.
\newenvironment{resumcap}{%
  \begin{tcolorbox}[blocdidactic]%
  {\color{DidacticVerd}\ifpdftex\scshape\else\sffamily\bfseries\fi S\'{\i}ntesi del cap\'{\i}tol.}\hspace{0.6em}%
  \ignorespaces
}{\end{tcolorbox}}